\runninghead{Scope and References}
\section{Appendix: Scope and Qualifications}
\label{sec:scope}

{\small
\textbf{Edition, Mass and occurrence.} The study treats the Mass
\latin{Dominica quinta quae superfuit post Epiphaniam}, II classis, of the
\work{Missale Romanum}, editio typica, Typis Polyglottis Vaticanis 1962,
printed pp.~420--421, nos.~1749--1758, as said on Sunday 8~November 2026, the
twenty-fourth Sunday after Pentecost, in the universal 1962 calendar and no
particular calendar. Its orations and lessons are those of the Fifth Sunday
after the Epiphany (pp.~47--48, nos.~326--335) and its chants those of the
Twenty-third Sunday after Pentecost (pp.~416--417), under RGMR~298. Its place
in 2026, the omission of the calendar commemoration of the Four Crowned
Martyrs, the colour, the Gloria, the Credo and the Trinity Preface rest on
the Missal's general rubrics (RG~10, 12, 14, 18~b, 19, 91, 107--111, 127~b;
RGMR~298, 431, 434, 475, 494~b) and on the Missal's printed direction after
the Secret, as \texttt{research/context.md} records; two dated Ordos of
institutes using the 1962 books, of the Fraternity of St Peter and of the
Institute of Christ the King Sovereign Priest, both in France, give the same
Mass for the date. No local oration, external solemnity or votive Mass is
considered. The \latin{Asperges}, the Ordinary and the Office are not treated.

\textbf{Text control.} Every Latin form was read on page images of the CMAA
facsimile of the 1962 typical edition and collated word for word with the
Ratisbon printing of 1862 (Pustet), whose page images carry every word of the
ten texts; the collation and the comparison with the Clementine Vulgate are in
\texttt{propers/verified.md}. Where the chants differ from the Clementine
Vulgate the differences are reported as textual facts; the Latin psalter or
antiphonary from which the chant forms came was not identified. The rubrics
of the 1962 book are described and not reproduced.

\textbf{English.} The English of the lessons, and of each chant's canonical
verses, is the Douay--Rheims in Challoner's revision. The English of the five
chants as the Missal sings them and of the three orations is that of the
Lasance \work{New Roman Missal} in the Benziger revision of 1945, public domain
in the United States by non-renewal, which prints this Mass as its
\work{Supplementary Mass No.~2}; its Postcommunion is quoted from p.~198,
which carries \latin{per haec mysteria}. None of this English is an approved
liturgical translation of the 1962 Missal. Augustine's sermons, letter and
psalm expositions, Chrysostom's homilies on Matthew and Colossians, and the
\work{Summa} are quoted in the nineteenth- and twentieth-century English
translations named in the References; Bellarmine in O'Sullivan's abridged
English; Schuster in the English of 1927. Elsewhere the Latin of the Fathers
is quoted with an English rendering made for this study, and no English is
composed for a liturgical text.

\textbf{Reception.} The interpretations rest on the loci named in the text and
recorded, with editions and bounds, in \texttt{research/scope.md} and
\texttt{research/interpretations.md}. Every connection drawn between two texts
of the Mass belongs to the study's interpretations unless an author is named as
making it. The Latin of Jerome, Hilary, Rabanus and Ambrose was read on page
images; Aquinas, Cassiodorus, Honorius, Rupert, Durandus, Sicard and Schuster
in uncorrected optical text or locally parsed web text; Chrysostom on the
Psalms, the Athanasian expositions, Theodoret and Theophylact in uncorrected
Greek or Latin optical text, so that Chrysostom on Psalms~43 and 129 is given
in paraphrase; Augustine on Psalm~129 in the abridged English of NPNF. No
column is given for Cassiodorus, the Athanasian expositions, Theodoret or
Honorius, whose optical text was read without page images; they are cited by
the verse or sermon they comment on. Theodoret's psalm
commentary, Origen on Matthew and Augustine's \work{Contra epistulam
Parmeniani} were not read. Every principal authority of an interpretation has
a row in the repository's register of authorial standing; Honorius,
Cassiodorus, Theodoret, Theophylact, Rupert, Durandus and Sicard are
ecclesiastical writers and are cited in support only, and the Psalm
expositions printed under St Athanasius's name are cited under that
attribution. Rupert, Durandus and Sicard expound these chants in Masses whose
lessons differed from this Mass's, and they are cited for the chants alone.
St Thomas Aquinas's conclusion on the punishment of heretics is reported as
his, and Augustine's \work{Contra epistulam Parmeniani} only as Aquinas cites
it; the study does not treat the later development of the Church's teaching
on religious liberty.

\textbf{Chronology.} Every biblical date in the Scriptural Date and Location
appendix is generated from the repository's Scripture chronology corpus, with
its sources named there. The record holds no date for the occasion of
Psalms~43, 84 or 129 or for the telling of the parable.

\textbf{Rights.} The Latin is published on 17 U.S.C. 103(b) together with its
public-domain witness, the Pustet printing of 1862. The Douay--Rheims and the
Lasance English are quoted from registered editions. The other quotations are
short and come from editions published before 1931 or from medieval and
ancient Latin: Migne's \work{Patrologia Latina}, vols.~9, 14, 24, 26, 70, 107,
170, 172 and 213, and \work{Patrologia Graeca}, vols.~27, 55, 82 and 124; the
Venice (1745), Parma (1863) and Dessain (1857) Aquinas; the Lyon Durandus
(1612); the Nicene and Post-Nicene Fathers; and Schuster's \work{Sacramentary}
(1927). The Corpus Thomisticum and Wikisource texts are used only for Latin
that is itself medieval or ancient. Scripture, liturgical texts and quoted
translations are not relicensed by the repository's licence.

\textbf{Use.} This is a study aid to the 1962 Missal, not an official
liturgical book, translation or critical edition.
\par}

\section{References}
\label{sec:references}

{\small
\subsection*{Liturgical books}
\begin{itemize}[leftmargin=1.2em,itemsep=.25em]
\item \work{Missale Romanum ex decreto Sacrosancti Concilii Tridentini
restitutum, Summorum Pontificum cura recognitum}, editio typica (Typis
Polyglottis Vaticanis, 1962), CMAA facsimile: \work{Rubricae generales} and
\work{Rubricae generales Missalis romani}; Proprium de Tempore, pp.~47--48,
nos.~326--335, and pp.~420--421, nos.~1749--1758.
\item \work{Missale Romanum} (Ratisbon: Pustet, 1862), pp.~42--44 and
357--359.
\item Lasance, F.~X., and Walsh, F.~A., \work{The New Roman Missal in Latin and
English}, revised edition (New York: Benziger Brothers, 1945), pp.~196--198
and 744--747.
\item Fraternité Saint-Pierre (France), \work{Ordo du mois}, and Institut du
Christ Roi Souverain Prêtre (France), \work{Ordo}, entries for 8~November 2026
(websites), read 5~November 2026.
\end{itemize}

\subsection*{Scripture}
\begin{itemize}[leftmargin=1.2em,itemsep=.25em]
\item \work{The Holy Bible Translated from the Latin Vulgate} (Douay--Rheims,
Challoner revision), Project Gutenberg text: Jer 29:1--14; Ps 43:1--10, 22;
84:1--5; 129:1--4; Mt 13:1--3, 24--43; Mk 11:20--25; Col 3:1--17.
\item \work{Biblia Sacra Vulgatae Editionis} (Clementine Vulgate), for the
comparisons with the Missal's forms.
\end{itemize}

\subsection*{Fathers, Doctors and other authors}
\begin{itemize}[leftmargin=1.2em,itemsep=.25em]
\item Ambrose, \work{Enarrationes in XII Psalmos}, on Ps 43, nn.~21--22, PL~14
(Paris, 1845), col.~1154.
\item Athanasius (attributed), \work{Expositiones in Psalmos}, on Ps 84, the
opening summary and v.~2, PG~27.
\item Augustine, \work{Sermo} 73 (Sermon XXIII in NPNF), first series,
vol.~6; \work{Epistula} 93 to Vincentius, NPNF, first series, vol.~1;
\work{Enarrationes in Psalmos} 42, 43, 64, 84 and 129, NPNF, first series,
vol.~8; \work{Quaestiones Evangeliorum} I.10, Migne Latin text (Wikisource).
\item Bede, \work{In Marci evangelium expositio} III, on 11:22--25, Migne
Latin text (Wikisource).
\item Bellarmine, Robert, \work{Explanatio in Psalmos}, in John O'Sullivan's
English, \work{A Commentary on the Book of Psalms}, on Pss 43, 84 and 129.
\item Cassiodorus, \work{Expositio Psalmorum}, on Ps 43:8 and Ps 84:2, PL~70.
\item Chrysostom, John, \work{Homilies on Matthew} 46 and 67, NPNF, first
series, vol.~10; \work{Homilies on Colossians} 8 and 9, NPNF, first series,
vol.~13; \work{Expositio in Psalmos}, on Pss 43 and 129, PG~55.
\item Hilary of Poitiers, \work{Tractatus super Psalmos}, on Ps 129, nn.~2--3,
PL~9 (Paris, 1844), cols.~719--720; \work{In Evangelium Matthaei
commentarius} 13.1, PL~9, col.~993.
\item Honorius Augustodunensis, \work{Speculum Ecclesiae}, the sermon on the
cockle headed \latin{Dominica tercia}, PL~172.
\item Jerome, \work{Commentarii in Matthaeum} II, on 13:24--30 and 13:36--43,
PL~26 (Paris, 1845), cols.~89, 93--94; \work{In Hieremiam} V, on 29:5--14,
PL~24 (Paris, 1865), cols.~891--894.
\item Rabanus Maurus, \work{Commentaria in Matthaeum} IV, on 13:24--30, PL~107
(Paris, 1864), cols.~947--948.
\item Theodoret of Cyrus, \work{Interpretatio epistolae ad Colossenses}, on
3:15--16, PG~82; Theophylact of Ohrid, \work{Expositio in epistolam ad
Colossenses}, on 3:15--16, PG~124, Latin version, cols.~1261--1264.
\item Thomas Aquinas, \work{Super Evangelium S.~Matthaei lectura}, cap.~13, in
\work{Opera omnia}, vol.~3 (Venice: Bettinelli, 1745), pp.~183--185;
\work{Super Psalmos}, on Ps 43, nn.~1--4, Corpus Thomisticum text and
\work{Opera omnia}, vol.~14 (Parma, 1863); \work{Super epistolam ad
Colossenses}, cap.~3, lect.~3, Corpus Thomisticum text and \work{In omnes
D.~Pauli Apostoli epistolas commentaria}, vol.~2 (Liège: Dessain, 1857);
\work{Summa theologiae} II-II q.~10 a.~8 and q.~11 a.~3, in the English
Dominican translation.
\end{itemize}

\subsection*{Commentators on the Mass}
\begin{itemize}[leftmargin=1.2em,itemsep=.25em]
\item Durandus, William, \work{Rationale divinorum officiorum} VI.140, in the
Lyon printing of 1612, tomus secundus, p.~412.
\item Rupert of Deutz, \work{De divinis officiis} XII.23, PL~170,
cols.~329--330.
\item Schuster, Ildefonso, \work{The Sacramentary (Liber Sacramentorum)},
vol.~III (London: Burns Oates \& Washbourne, 1927), pp.~186--188.
\item Sicard of Cremona, \work{Mitrale} VI.23, PL~213.
\end{itemize}
\par}
